\documentclass[a4paper,11pt]{ltjsarticle}
\usepackage[margin=22mm]{geometry}
\usepackage{amsmath}
\usepackage{booktabs}
\usepackage{listings}
\usepackage{xcolor}
\usepackage{hyperref}
\hypersetup{colorlinks=true,linkcolor=blue,urlcolor=blue,citecolor=blue}

\lstset{
  basicstyle=\ttfamily\footnotesize,
  backgroundcolor=\color{black!4},
  frame=single,
  framerule=0pt,
  breaklines=true,
  columns=fullflexible,
}

\title{Embedded Xinu (Raspberry Pi 3) 上の\\GC 連動・分散 N-Queens ベンチマーク報告}
\author{Yasushi Kodama}
\date{2026年6月2日}

\begin{document}
\maketitle

\section{概要}
ベアメタル Embedded Xinu (arm-rpi3 / BCM2837 / Cortex-A53) 上で動作する
AIPL アクター・ランタイムの負荷分散フレームワーク (Dispatcher + Worker 群) を用い、
\textbf{周期 GC アクターを有効にした状態で分散 N-Queens 問題のベンチマーク}を実施した。
解の正しさ・スループット・GC の動作を確認し、計測中に判明した制御プレーンの
ボトルネックと、その改善 (Collector への push 集約) を報告する。
ビルド ID は \texttt{xinu-gwm-b69}。

\section{システム構成}
\begin{itemize}
  \item \textbf{Pi 3 (arm-rpi3)}: 有線 Ethernet 経由の HTTP サーバ \texttt{webactor}
        (192.168.3.50:8080) が AIPL ランタイムを制御。
  \item \textbf{負荷分散}: \texttt{Dispatcher} 1 体 + \texttt{Worker} 4 体
        (\texttt{/api/loadbal/init})。Dispatcher は最小負荷のワーカーへタスクを振り分ける。
  \item \textbf{GC}: \texttt{Collector} アクター (\texttt{/api/gc-actor/init},
        \texttt{enable}) を起動。周期 5\,s でアクター表をスイープし、
        アイドル閾値 30\,s を超えた使い捨てアクターを回収する。
  \item \textbf{計測ドライバ}: Mac 側 \texttt{tools/nqueens\_bench.py}
        (参照解は Mac で計算し照合)。
\end{itemize}

\section{ベンチマーク手法}
$N$-Queens を\textbf{第 1 列ごとに分割}し、各列を 1 個の \texttt{compute\_nq}
タスクとして Dispatcher へ投入する (\texttt{/api/loadbal/nqueens?n=N})。
各 Worker が担当列の部分解数を数え、Dispatcher へ \texttt{done(idx, task\_id, result)}
を返す。全列の部分解数の総和が $N$-Queens の総解数となる。
GC を有効にしたまま実行し、計算アクターの生成・回収が並行して行われる状況を再現した。

\section{結果}
\subsection{解の正しさとスループット}
\begin{table}[h]\centering
\begin{tabular}{cccc}
\toprule
$N$ & 分散実行 (Pi 3) 総解数 & 既知の正解 & 判定 \\
\midrule
8 & \textbf{92}  & 92  & 一致 \\
9 & \textbf{352} & 352 & 一致 \\
\bottomrule
\end{tabular}
\caption{分散 N-Queens の総解数 (Pi 3、列分散・GC 有効)}
\end{table}

$N=8$ の第 1 列ごとの内訳は左右対称で
$\{4,8,16,18,18,16,8,4\}$、総和 92 と一致した。
各 \texttt{compute\_nq} タスクの実計算時間は Pi 3 上で約 \textbf{20\,ms/タスク}。
照合用の Mac 参照計算は $N=8$ で 92 解を 2\,ms。

\subsection{GC の動作}
ベンチ実行中の GC アクター統計を以下に示す。
\begin{lstlisting}
gc-actor obj=6 enabled=1  period_ms=5000 threshold_ms=30000
sweep_count=1  swept_total=0  last_scanned=7
\end{lstlisting}
周期スイープが実行され (\texttt{sweep\_count}$\ge$1)、アクター 7 体をスキャンした。
\texttt{swept\_total=0} は\textbf{回収対象なし}を意味する。Dispatcher・Worker・Collector
は常駐であり、N-Queens タスクは部分解数を返すのみで使い捨てアクターを残さないため、
回収すべきゴミが発生しないのは正しい挙動である。

\section{考察: 制御プレーンのボトルネック}
分散計算そのものは正しく高速 (約 20\,ms/タスク) であったが、
\textbf{Pi 3 の \texttt{webactor} は単一スレッド}であり、
Worker の N-Queens 計算が走ると HTTP 制御・集計が一時的に CPU 飢餓になる。
当初の集計方式は Mac が各タスクを \texttt{/api/loadbal/task?id=K} で個別ポーリングして
合計する方式であったため、計算負荷下で HTTP 応答が滞り、
全自動ベンチ (\texttt{nqueens\_bench.py}) が 300\,s の予算内に集計を完了できず
タイムアウトした。すなわちボトルネックは\textbf{分散計算ではなく、
計算と同一コアを奪い合う制御プレーン (HTTP ポーリング)}にある。

\section{改善: Collector への push 集約}
上記を受け、集計を\textbf{ポーリングから push 集約へ}変更した (ビルド \texttt{b69})。
\begin{itemize}
  \item Worker は従来どおり完了時に Dispatcher へ \texttt{done} を\textbf{push}する。
  \item Dispatcher の \texttt{done} ハンドラで N-Queens タスク (kind=\texttt{'n'}) の
        部分解数をその場で\textbf{積算}する (\texttt{g\_nq\_sum},
        \texttt{g\_nq\_received})。
  \item Mac は新設の \texttt{/api/loadbal/nqueens-result} を\textbf{1 回}問い合わせるだけで、
        集約済みの総和と完了状況 (\texttt{done/received/expected/total}) を得る。
\end{itemize}
これにより、$N$ 個のタスクに対する HTTP 集計リクエスト数が
\textbf{$O(N)$ 回 $\rightarrow$ 1 回}へ削減され、計算と制御プレーンの競合が大幅に緩和される。
集約は Pi 3 側 (Dispatcher) で完結し、Mac 側は受動的に総和を読むだけとなる。

\section{結論}
GC を有効にした状態で、分散 N-Queens は $N=8 \to 92$、$N=9 \to 352$ と
\textbf{正答}を得た。GC は周期スイープを実行し (常駐アクターのみのため回収 0)、
計算と並行して正常に動作した。スループット上の制約は分散計算ではなく単一スレッド
HTTP 制御プレーンにあることを特定し、\textbf{Collector への push 集約}
($O(N)\to1$ の集計リクエスト削減) として改善を実装した。

\end{document}
